\documentclass[10pt]{article}
\usepackage[margin=12mm]{geometry}
\usepackage{graphicx}
\usepackage{caption}
\pagestyle{empty}
\captionsetup{font=small,labelfont=bf}

\begin{document}
\begin{figure}[p]
\centering
\includegraphics[width=\textwidth]{terminal_main_pareto_front.pdf}
\caption{\textbf{The terminal layer of the \textit{C. elegans} lineage is
nearly Pareto optimal.} Terminal-cell identities, final positions, and
protein-expression states were held fixed while their direct parentage edges
were reconstructed using optimal linear assignment under varying distance
objectives defined in Equation \ref{eqn:pareto}. The resulting curve is the
exact Pareto front for travel distance and cell-state distance. Distances are
measured in first-cousin-null standard deviations and translated so the
natural lineage (black cross) is at the origin. Relaxing the lineage
neighborhood within which assignments can be shuffled moves the null
distributions progressively towards the upper-left: first-cousin (yellow),
second-cousin (orange), third-cousin (green), and full-random shuffling
(magenta inset). The outlined circle marks the Pareto-optimal assignment with
maximum edge retention compared with natural lineage, and is filled with
its corresponding color from the edge-retention scale with which the assignments
along the Pareto front are colored by edge retention; brighter blue
denotes greater retention.}
\label{fig:ce-terminal-pareto-main}
\end{figure}
\end{document}
